%  ----------------------------------------------------------------------------
%
%       Copyright (for the thesis) 2009 by [author - insert yourself]
%
%       This thesis is published under the
%       Creative Commons Attribution-No Derivative Works 3.0 Austria License
%       as detailed at http://creativecommons.org/licenses/by-nd/3.0/at/
%
%  ----------------------------------------------------------------------------
%  Template credits and license:
%  ----------------------------------------------------------------------------
%
%       "Fakult�t f�r Informatik" diploma/master thesis template 2008
%
%       based upon "Diploma thesis template 2005" by lukas.silberbauer(at)gmx.at
%       based upon "Diplomarbeit mit LaTeX" by Tobias Erbsland
%       incorporating a title page by Informatik-Forum user "Baby"
%       polished and ported to the TU fonts package by Jakob Petsovits
%
%       published under the terms of
%
%  ----------------------------------------------------------------------------
%  "THE BEER-WARE LICENSE":
%  <lukas.silberbauer(at)gmx.at> wrote this file. As long as you retain this
%  notice you can do whatever you want with this stuff. If we meet some day,
%  and you think this stuff is worth it, you can buy me (us) a beer in return.
%  ----------------------------------------------------------------------------
%
%  (end of template credits)
%

%
% Variablen, die f�rs Deckblatt verwendet werden.
% Die grobe Dokumentenstruktur findet sich übrigens ganz unten in dieser Datei.
%
\newcommand{\thesistype}{TUTORIAL}
\newcommand{\thesistitle}{Scalimero}
\newcommand{\thesissubtitle}{ }

\newcommand{\thesisauthor}{Gabriel A. Grill and Alexander C. Steiner}
\newcommand{\matrikelnr}{ }
\newcommand{\zuerlangendertitel}{ }
\newcommand{\studium}{ }

\newcommand{\institut}{Institute of Computer Aided Automation}
\newcommand{\institutsuni}{Vienna University of Technology} % disabled by default
\newcommand{\betreuer}{Ao.Univ.Prof. Dr. Wolfgang Kastner}
\newcommand{\assistent}{ }

%
% Globale Weichenstellungen f�rs ganze Dokument.
%
\documentclass[%
  pdftex,%              PDFTex verwenden
  a4paper,%             A4 Papier
  oneside,%             Einseitig
  bibtotocnumbered,%    Literaturverzeichnis nummeriert einfügen
  idxtotoc,%            Index ins Verzeichnis einfügen
  halfparskip,%         Europäischer Satz mit abstand zwischen Absätzen
  chapterprefix,%       Kapitel anschreiben als Kapitel
  %headsepline,%         Linie nach Kopfzeile
  %footsepline,%         Linie vor Fusszeile
  12pt%                 Größere Schrift, besser lesbar am Bildschrim
]{scrbook}


%
% Paket für die Indexerstellung.
%
\usepackage{makeidx}

%
% Paket für Übersetzungen ins Deutsche
%
\usepackage[german, english]{babel}

%
% Paket, um \today auch im Format dd.mm.yyyy anzeigen zu können
%  (für die Titelseite, mit \numdate)
%
\usepackage[german]{isodate}

%
% Pakete um Latin1 Zeichnensätze verwenden zu können und die dazu
%  passenden Schriften.
%
\usepackage[utf8]{inputenc}
\usepackage[T1]{fontenc}

%
% Paket zum Erweitern der Tabelleneigenschaften
%
\usepackage{array}

%
% Paket um Grafiken einbetten zu können
%
\usepackage{graphicx}


%
% Zeilenabstand einstellen
%
\usepackage{setspace}
\onehalfspacing
%\doublespacing


%\setlength{\baselineskip}{24pt}
%\renewcommand{\baselinestretch}{1.5}


%
% define variables
%
\def\maintitle#1{\gdef\maintitle{#1}}
\def\subtitle#1{\gdef\subtitle{#1}}

%
% Zeilenumbruch bei Bildbeschreibungen.
%
\setcapindent{1em}

%
% kopf und fusszeilen
%
\pagestyle{headings}

%
% mathematische symbole aus dem AMS Paket.
%
\usepackage{amsmath}
\usepackage{amssymb}

%
% Type 1 Fonts für bessere darstellung in PDF verwenden.
%
\usepackage{mathptmx}           % Times + passende Mathefonts
%\usepackage[scaled=.92]{helvet} % skalierte Helvetica als \sfdefault
\usepackage{courier}            % Courier als \ttdefault
\usepackage{tufonts}            % TU-Schriften als \sfdefault

%
% Spezielle Schrift verwenden.
%
\renewcommand{\encodingdefault}{T1}
\renewcommand{\familydefault}{\rmdefault}
\renewcommand{\bfdefault}{db}   % bold would be too bold, use demibold instead
\renewcommand{\mddefault}{l}    % use TU Light, looks better than medium weight



%
% Paket um Textteile drehen zu können
%
\usepackage{rotating}


%
% Für Acronyme
%
\usepackage{acronym}

%
% Package für Farben im PDF
%
\usepackage{color}

%
% Paket für Links innerhalb des PDF Dokuments
%
\definecolor{LinkColor}{rgb}{0,0,0.5}

\usepackage[
pdfauthor={\thesisauthor},
bookmarks=true, % PDF bookmarks allowed. NB! The level depth of bookmarks is the same as in the TOC.
unicode=true, % PDF bookmarks in Unicode.
bookmarksnumbered=true, % Section numbers in PDF bookmarks.
bookmarksopenlevel=1, % The open level in PDF bookmarks.
hyperindex=true, % Hyperlinked index.
plainpages=false, % Name arabic and roman page numbers differently.
colorlinks=true, % Links are marked as coloured text, not coloured box.
linkcolor=linkc, % The colour for in-document links (e.g. in the table of contents).
citecolor = citec, % The colour for bibliographic citations.
urlcolor=urlc, % The colour for hyperlinks to the Net.
pdfpagelayout=OneColumn % Continuous page scrolling.
]{hyperref}
\hypersetup{
  colorlinks=true,
  linkcolor=LinkColor,
  citecolor=LinkColor,
  filecolor=LinkColor,
  menucolor=LinkColor,
  pagecolor=LinkColor,
  urlcolor=LinkColor
}

%
% Paket um Listings sauber zu formatieren.
%
\usepackage[savemem]{listings}
\lstloadlanguages{TeX}

%
% ---------------------------------------------------------------------------
% Listing Definitionen für PHP Code
%
\definecolor{lbcolor}{rgb}{0.85,0.85,0.85}
\lstset{language=[LaTeX]TeX,
  numbers=left,
  stepnumber=1,
  numbersep=5pt,
  numberstyle=\tiny,
  breaklines=true,
  breakautoindent=true,
  postbreak=\space,
  tabsize=2,
  basicstyle=\ttfamily\footnotesize,
  showspaces=false,
  showstringspaces=false,
  extendedchars=true,
  backgroundcolor=\color{lbcolor}
}

% ---------------------------------------------------------------------------
% Neue Umgebungen
% ---------------------------------------------------------------------------

\newenvironment{ListChanges}%
  {\begin{list}{$\diamondsuit$}{}}%
  {\end{list}}

%
% Index erzeugen
%
\makeindex


%%%%%%%%%%%%%%%%%%%%%%%%%%%%%%%%%%%%%%%%%%%%%%%%%%%%%%%%%%%%%%%%%%%%%%
%
% Document structure
%
%%%%%%%%%%%%%%%%%%%%%%%%%%%%%%%%%%%%%%%%%%%%%%%%%%%%%%%%%%%%%%%%%%%%%%

\begin{document}

%\pagenumbering{alph}
%  ----------------------------------------------------------------------------
%
%       Copyright (for the thesis) 2009 by [author - insert yourself]
%
%       This thesis is published under the
%       Creative Commons Attribution-No Derivative Works 3.0 Austria License
%       as detailed at http://creativecommons.org/licenses/by-nd/3.0/at/
%
%  ----------------------------------------------------------------------------
%  Template credits and license:
%  ----------------------------------------------------------------------------
%
%       "Fakult�t f�r Informatik" diploma/master thesis template 2008
%
%       based upon "Diploma thesis template 2005" by lukas.silberbauer(at)gmx.at
%       based upon "Diplomarbeit mit LaTeX" by Tobias Erbsland
%       incorporating a title page by Informatik-Forum user "Baby"
%       polished and ported to the TU fonts package by Jakob Petsovits
%
%       published under the terms of
%
%  ----------------------------------------------------------------------------
%  "THE BEER-WARE LICENSE":
%  <lukas.silberbauer(at)gmx.at> wrote this file. As long as you retain this
%  notice you can do whatever you want with this stuff. If we meet some day,
%  and you think this stuff is worth it, you can buy me (us) a beer in return.
%  ----------------------------------------------------------------------------
%
%  (end of template credits)
%

\selectlanguage{german}
\begin{titlepage}
\fontfamily{fts}\selectfont

\vspace*{-0.75cm}

\begin{minipage}[t]{0.3\linewidth}
\begin{flushleft}
\includegraphics[width=3.49cm,height=1.65cm]{images/fakultaet-fuer-informatik-logo.png}
\end{flushleft}
\end{minipage}
\begin{minipage}[t]{0.7\linewidth}
\begin{flushright}
\vspace*{-1.65cm}
\includegraphics[width=9.34cm,height=0.53cm]{images/fakultaet-fuer-informatik-text.png}
\end{flushright}
\end{minipage}

\vspace{1.0cm}

\begin{center}
\dmseries\huge{\thesistitle}
\end{center}

\vspace{0.5cm}

\renewcommand{\bigskip}{\vspace*{0.43cm}}

\begin{center}
\begin{Large}\thesistype \end{Large} \bigskip \\
%\begin{normalsize}zur Erlangung des akademischen Grades \end{normalsize} \bigskip \\
%\begin{Large}\dmseries\zuerlangendertitel\end{Large} \bigskip \\
%\begin{normalsize}im Rahmen des Studiums\end{normalsize} \bigskip \\
%\begin{large}\dmseries\studium\end{large} \bigskip \\
\begin{normalsize}ausgef\"{u}hrt von \end{normalsize} \bigskip \\
\begin{large}\dmseries\thesisauthor\end{large} \\
%\begin{normalsize}Matrikelnummer \matrikelnr\end{normalsize} \\
\end{center}

\vspace{2cm}

\begin{flushleft}
am: \\ \textit{\institut} \\
%\small{der \institutsuni} \\
\end{flushleft}


\begin{flushleft}
\smallskip
%Betreuung: \\
\textit{Betreuer: \betreuer} \\
%\textit{Mitwirkung: \assistent}
\end{flushleft}

\vspace{3cm}

\date{\today}
\numdate

\footnotesize
\vspace{1.5cm}

\begin{flushleft}
\begin{tabbing}%
%Wien, \today & \hspace*{0.5cm} & \underline{\hspace*{5.2cm}} & \hspace*{1.0cm} & \underline{\hspace*{5.2cm}}\\
%                       &                 &         Unterschrift Verfasserin  & & Unterschrift Betreuer
\hspace{4.5cm} \= \hspace{5cm} \= \hspace{5cm} \kill
\textit{Wien, \today} \> \underline{\hspace*{5cm}} \> \hspace{0.8cm} \underline{\hspace*{5cm}} \\
     \> \hspace{0.1cm} (Unterschrift Verfasser/in) \> \hspace{1.0cm} (Unterschrift Betreuer/in) \\
\end{tabbing}
\end{flushleft}

\vspace{0.9cm}

\begin{minipage}[b]{1.0\linewidth}
\begin{center}
\underline{\hspace*{15.5cm}}
\scriptsize{
Technische Universit\"{a}t Wien \\
A-1040 Wien \hspace{0.5cm} Karlsplatz 13 \hspace{0.5cm} Tel. +43/(0)1/58801-0 \hspace{0.5cm} http://www.tuwien.ac.at
}
\end{center}
\end{minipage}

\vspace*{-2.5cm}

\end{titlepage}

%
% EOF
%


%\selectlanguage{english}
%\pagenumbering{roman}
%\setcounter{page}{1}

%\input{chapters/abstract}
%\input{chapters/acknowledgements}

\tableofcontents

%\pagenumbering{arabic}
%\setcounter{page}{1}

\chapter{Introduction}

Because of the increasing poularity of home automation, the desire to develop applications for controlling your house became more and more important. Calimero is a collection of Java APIs that together form a foundation for building such applications in EIB/KNX installations. Detailed knowledge of the protocol is not required. On the top of that Calimero only requires J2ME enviorments, which enables use on embedded platforms. It is an open source project and was developed by the Institute of Computer Aided Automation of the Technical University of Vienna.

EIBnet/IP allows you to communicate with components of the widely spreaded EIB/KNX standard by tunnelling over IP Networks. The protocol is published in the KNX Handbook  and in the European standard EN 13321-2:2006 (Open Data Communication in Building Automation, Controls and Building Management - Home and Building Electronic Systems - Part 2: KNXnet/IP Communication), available from European standardization organizations. Note: EN 13321-2 is useless without information on the KNX/EIB control network specific data structures (cEMI, DPTs); these should be defined in EN 50090 (Home and Building Electronic Systems). 

\clearpage
Since development of Calimero many great projects are using it:
\begin{itemize}
	\item \textbf{KNX@Home}(\url{http://knxathome.fh-deggendorf.de/})
	\item \textbf{BASys}(\url{http://sourceforge.net/projects/basys/})
	\item \textbf{KNXnet/IP Wireshark dissector}(\url{http://knxnetipdissect.sourceforge.net/})
	\item \textbf{CONECT}(\url{http://sourceforge.net/projects/conect/})
	\item \textbf{EIB Home Server}(\url{http://eibcontrol.sourceforge.net/})
	\item \textbf{LEIBnix}(\url{http://leibnix.sourceforge.net/Wikka/HomePage})
	\item \textbf{Sombrero}(\url{http://grill.github.com/sombrero/})
	\item \textbf{Scalimero}(\url{http://grill.github.com/SCalimero/})
\end{itemize}
\clearpage

\chapter{Installation}

First, download the calimero package from \url{http://www.auto.tuwien.ac.at/downloads/calimero-all-2.0a4.zip} and unpack it to a detination of your choice. It contains several other archives, including 4 Java Archives (.jar), the source code of all of them and documentation of tools and the library. The core library is calimero-2.0a4.jar, this is the file you want to have on the classpath of all your calimero projects. Calimero-gui and calimero-tools are graphical and command line tools to test calimero and KNX, and to gain a better understanding of calimero by looking at their source code. The calimero-rxtx package contains an optional way of serial port access, which is not needed in most circumstances.

The important thing is to put calimero-2.0a4.jar on your classpath. How to add files to your classpath depends on your development environment and should easily be found in the corresponding manual or home page. It is a good idea to unzip the calimero-2.0a3.zip to a place where you can access the documentation comfortably. It is also useful to get familiar with the command line tools and especially their source code, because it covers some important use cases in a concise way.

\chapter{DSL}
The scalimero DSL is an easy interface for accessing KNX devices and meant to be used in the Scala interpreter. It features convenient classes and a lot of implicit conversions to make the code look good while being able to check type bounds.

\section{Example}
To use the DSL,
\begin{lstlisting}
import tuwien.auto.scalimero.dsl._
\end{lstlisting}
Then create a network and open it using
\begin{lstlisting}
Network("10.0.0.5") open
\end{lstlisting}
The IP address specified here is the KNX router used to access the network. Then, we need to create some devices.
\begin{lstlisting}
val lA = Lamp("1/1/0")
val lB = Lamp("1/1/1")
\end{lstlisting}
Then we can turn the lamps on and off.
\begin{lstlisting}
lA turn on
lA turn off
\end{lstlisting}
More general:
\begin{lstlisting}
lB send true
lB send false
\end{lstlisting}
We can also read from the devices.
\begin{lstlisting}
val b = lA.read
\end{lstlisting}
Note that b is of type \lstinline!Boolean!, it gets converted from KNX data automatically.

You can also subscribe to events:
\begin{lstlisting}
lA.eventSubscribe(on) {
  println("lA has been turned on")
}
\end{lstlisting}

If you want to be notified for every write on the device, subscribe a write callback:
\begin{lstlisting}
lA.writeSubscribe{
  newstatus : Boolean =>
  println("lA status: " + newstatus)
}
\end{lstlisting}

\section{Network}

The \lstinline!Network! class and companion object in the \lstinline!tuwien.auto.scalimero.connection! package are the access point to the KNX network. To declare a \lstinline!Network!, use the companion objects apply method:

\begin{lstlisting}
Network("knxrouter")
\end{lstlisting}

Replace \lstinline!"knxrouter"! with a correct KNX router IP address or DNS name. The default medium is twisted pair 1 (9600 bit/s). If you need a different medium for a particular \lstinline!Network!, use the optional medium parameter:

\begin{lstlisting}
Network("knxrouter", TPSettings.TP0)
\end{lstlisting}

You can set this globally for all \lstinline!Network!s by assigning to the companion object's \lstinline!defaultMedium! field:

\begin{lstlisting}
Network.defaultMedium = TPSettings.TP0
Network("knxrouter")
\end{lstlisting}

All devices will use the last \lstinline!Network! created by default. If you want to create the \lstinline!Network!s seperately from the devices, you can pass a code block to the \lstinline!Network!s, and every device created within will use the specific \lstinline!Network!. It gets clearer in an example:

\begin{lstlisting}
val newnet = Network("knxrouter")
val oldnet = Network("knxrouter-old", TPSettings.TP0)

newnet {
  Lamp("1/1/1")
  Lamp("1/1/2")
}
oldnet {
  Lamp("1/2/3")
  Lamp("1/2/5")
}
\end{lstlisting}

While this was introduced to resemble a markup language, it may be difficult to access the devices from the outside.

If you want to update a particular device (and trigger its callbacks), you can send a read request to the KNX device. It will answer with its current value and the \lstinline!Network! will relais that update to the device objects:

\begin{lstlisting}
net.readRequest("1/1/1")
\end{lstlisting}

Note that this call is asynchronous and can (and should) be done from the device object, if available, using its \lstinline!readRequest! method.

You can also subscribe to events on the KNX network. This is seldom needed, however, as devices do so automatically and messages received from the \lstinline!Network! lack type information. Still, it can be useful if you need to implement a seperate device handling mechanism:

\begin{lstlisting}
val net = Network("knxrouter")
val act = actor {
  case e : WriteEvent => ...
  case e : ProcessEvent => ...
}
net subscribe act
\end{lstlisting}

The actor will start to receive \lstinline!tuwien.auto.scalimero.WriteEvent!s and \lstinline!tuwien.auto.calimero.process.ProcessEvent!s. \lstinline!WriteEvent!s occur when a scalimero device writes something to the network while \lstinline!ProcessEvent!s indicate that something on the network (e.g. a switch) sent an update. In most cases, however, it is easier to subscribe to devices themselves.

\section{Devices}

Devices are abstractions for the physical devices connected to the KXN bus. In most cases, the preconfigured devices in \lstinline!tuwien.auto.scalimero.device.preconf! should suffice (for information on how to create your own devices and advanced funcionality, see later chapters). As stated above, devices automatically subscribe to the last created \lstinline!Network!, so it is a bad idea to create a device without creating a \lstinline!Network! first. To create a device, use one of the constructor objects in \lstinline!preconf!:

\begin{lstlisting}
val lamp = Lamp("1/1/1")
val switch = Switch("1/1/2")
val dimmer = Dimmer("1/1/3")
val rollerBlind = RollerBlind("1/1/4")
val temperature = Temperature("1/1/5")
\end{lstlisting}

\lstinline!Lamp! and \lstinline!Switch! use boolean values or the more descriptive \lstinline!on! and \lstinline!off! aliases, while \lstinline!Dimmer! and \lstinline!RollerBlind! use percent values (= \lstinline!Int!) from 0 to 100 and \lstinline!Temperature! uses float values (unit depends on your thermostat).

To set values, use the send method, or \lstinline!Lamp!'s and \lstinline!Switch!'s more natural \lstinline!turn! method.

\begin{lstlisting}
lamp turn on
switch turn off
dimmer send 100
rollerBlind send 0
temperature send 19
\end{lstlisting}

To read values, use the read method. It always is of the appropriate type.

You can also subscribe callbacks. There are two types of callbacks, event callbacks and write callbacks. Event callbacks register to a specific type of event and execute whenever that event occurs while write callbacks get called every time the state of a device is updated.

\begin{lstlisting}
lamp.eventSubscribe(on) {
  println("lA has been turned on")
}

dimmer writeSubscribe {
  newVal : Int =>
  println("new dimmer value:" newVal)
}
\end{lstlisting}

The subscription methods return \lstinline!WriteCallback! or \lstinline!EventCallback! objects, which can be used to change the event code or detach(unsubscribe) the event.

\begin{lstlisting}
val wc = dimmer writeSubscribe {
  newVal : Int =>
  println("new dimmer value:" newVal)
}

def loggit(l : Logger) {
  wc update {
    newVal : Int =>
    l.info("new dimmer value:" newVal)
  }
}

def shutdown {
  wc.detach
  //which is the same as
  //dimmer writeUnsubscribe wc
}
\enc{lstlisting}

\chapter{Devices}
\section{Overview}

In Calimero KNX devices were represented through the \lstinline!Datapoint! class. Instances of this class could be used to send and receive messages by passing it as a parameter to a \lstinline!ProcessCommunicator!.

In Scalimero a more decentralized approach was chosen. Devices were introduced.

\paragraph{Devices have the following members (constructor parameters):}
\begin{itemize}
  \item \textbf{A KNX address} - Name: destAddr

    Can be passed as a \lstinline!String! or \lstinline!GroupAddress! object to the constructor of \lstinline!Device!.
  \item \textbf{DPType} - Name: dpt

    \lstinline!DPType! is a wrapper class for \lstinline!DPT!, which means a instance of \lstinline!DPType! or \lstinline!DPT! can be passed as a constructor parameter.
  \item \textbf{Name}[optional] - Name: name

    The name is a \lstinline!String! an will be given to the \lstinline!Datapoint! instance encapsulated in the \lstinline!Device! class.
  \item \textbf{Network}[optional] - Name: net

    If this parameter is not given, the value of \lstinline!Network.default!, which contains a reference to last created \lstinline!Network! instance, will be used.
\end{itemize}

\paragraph{Device traits:}
\begin{itemize}
  \item \textbf{TDevice} - Every device inherits this trait.
  \item \textbf{TCommandDevice} - Every device with the \lstinline!send! method inherits this trait.
  \item \textbf{TStateDevice} - Every device with the \lstinline!read! method inherits this trait.
\end{itemize}

The \lstinline!Device! base class is an \lstinline!Actor!, which means it can send and receive messages. Because of this ability it is possible to send \lstinline!ProcessEvent!s or \lstinline!WriteEvent!s to the device actor to trigger subscribed events.

\paragraph{Subclasses of Device:}
\begin{itemize}
  \item \textbf{CommandDevice}

  Devices of this class have no state, which means they can only be used for sending messages.

    \textbf{send}(d: DataPointValue)

        \lstinline!DataPointValue! subclasses are wrapper classes for primitive types. They are specified in the \lstinline!tuwien.auto.scalimero.device.dtype! package. It is possible to pass an instance of \lstinline!DataPointValue! or a primitive type value.
  \item \textbf{StateDevice}

  This class has the ability to send and read messages. Even though the name of the class is \lstinline!StateDevice!, no state is stored in the objects. It has got two methods to read values from the KNX Network:
    \begin{itemize}
        \item \textbf{read}: PrimitveType

            Returns a value with the primitive type chosen through the \lstinline!DPType!. If network communication fails, an Exception is thrown.
        \item \textbf{readOption}: Option[PrimitiveType]

            Returns a value with the primitive type chosen through the \lstinline!DPType! encapsulated in an Option.
    \end{itemize}
\end{itemize}

\section{Device in the REPL}
It is advised to use the base Device classes only for the purpose of testing or quick scripting. In most cases it is more rewarding to create a specialzied Device class and let it inherit from \lstinline!StateDevice! or \lstinline!CommandDevice!. How this can be done exactly will be explained in the next chapter.

To use it,
\begin{lstlisting}
import tuwien.auto.scalimero._
import device._
import device.dtype._
import connection._
\end{lstlisting}
Then create a network and open it. Now the creation of the devices can begin.
\begin{lstlisting}
val d1 = Device("1/1/1", Boolean.SWITCH)
\end{lstlisting}
To shorten the qualification of the \lstinline!DPType! just
\begin{lstlisting}
import device.dtype.Boolean._
\end{lstlisting}
 Now all Boolean \lstinline!DPType!s are in scope and it is possible to write
\begin{lstlisting}
val d2 = Device("1/1/0", SWITCH)
\end{lstlisting}
By writing the above code the \lstinline!apply! method of the companion object \lstinline!Device! is called and instantiates a \lstinline!StateDevice! with the given parameters.

To send messages
\begin{lstlisting}
d1 send true
\end{lstlisting}
or
\begin{lstlisting}
d2 send on
\end{lstlisting}
To read messages
\begin{lstlisting}
try{
    if(d1 read)
        println("Lamp was turned on")
    else
        println("Lamp was turned off")
} catch {
    case e => println("No value was received, because an error occured") }
\end{lstlisting}
or
\begin{lstlisting}
d1 readOption match {
    case Some(value) =>
        if(value)
            println("Lamp was turned on")
        else
            println("Lamp was turned off")
    case _ => println("No value was received, because an error occured")
}
\end{lstlisting}
The \lstinline!Device! base classes can only subscribe to \lstinline!WriteCallback!s.
\section{SimpleDevice}

The class \lstinline!SimpleDevice! was created because of the need for a quick to write abstraction with little overhead. This is accomplished by removing the Actor functionality. This class has the same \lstinline!read! and \lstinline!send! methods as \lstinline!StateDevice! and some additional low level sending methods:

\begin{itemize}
  \item \textbf{readRequest} - sends a read request
  \item \textbf{write}\lstinline!(value: PrimitiveType)! - writes a value to the KNX device
  \item \textbf{write}\lstinline!(value: Array[Byte])! - writes a value to the KNX device
  \item \textbf{write}\lstinline!(value: String)! - writes a value to the KNX device (string format depends on \lstinline!DPT!)
\end{itemize}
\clearpage
\section{Grouping of Devices}

Grouping of Devices can be done on different levels:

\begin{itemize}
  \item \textbf{low level} - through \lstinline!WriteCallback!s
  \item \textbf{mid level} - through events
  \item \textbf{high level} - through:
    \begin{itemize}
        \item \textbf{GroupDevice}
        
        \lstinline!GroupDevice! inherits from \lstinline!HashSet! and because of that it can be filled with \lstinline!Device!s. With the \lstinline!send! method a message can be sent to all \lstinline!Device!s in the \lstinline!GroupDevice!. The \lstinline!setMaster! method enables you to set a master device. When the state of the master device changes, all \lstinline!Device!s in the \lstinline!GroupDevice! are set to the same state. Through the \lstinline!addProxyFunction! method it is possible set a function which will be evaluated before the new received state is forwarded to the other devices. The result of this function will then be sent instead of the original received value.
        

\begin{lstlisting}
val gd = GroupDevice(d1, d2)

//all lights on
gd send on
//or
gd map {_ send on}

gd setMaster d1

//d2 will also be turned off automatically
d1 send off
\end{lstlisting}
        \item \textbf{MultipleAddressDevice}

        This class is an abstraction for devices with a separate reading and writing address.
\begin{lstlisting}
val mad1 = new MultipleAddressDevice(d1, d2)
//or
val mad2 = new MultipleAddressDevice("1/1/1", "1/1/0", Switch)
\end{lstlisting}
    \end{itemize}
\end{itemize}

\chapter{How to develop Devices}
\section{Getting started}
\section{Events}
%  ----------------------------------------------------------------------------
%
%       Copyright (for the thesis) 2009 by [author - insert yourself]
%
%       This thesis is published under the
%       Creative Commons Attribution-No Derivative Works 3.0 Austria License
%       as detailed at http://creativecommons.org/licenses/by-nd/3.0/at/
%
%  ----------------------------------------------------------------------------
%  Template credits and license:
%  ----------------------------------------------------------------------------
%
%       "Fakultät für Informatik" diploma/master thesis template 2008
%
%       based upon "Diploma thesis template 2005" by lukas.silberbauer(at)gmx.at
%       based upon "Diplomarbeit mit LaTeX" by Tobias Erbsland
%       incorporating a title page by Informatik-Forum user "Baby"
%       polished and ported to the TU fonts package by Jakob Petsovits
%
%       published under the terms of
%
%  ----------------------------------------------------------------------------
%  "THE BEER-WARE LICENSE":
%  <lukas.silberbauer(at)gmx.at> wrote this file. As long as you retain this
%  notice you can do whatever you want with this stuff. If we meet some day,
%  and you think this stuff is worth it, you can buy me (us) a beer in return.
%  ----------------------------------------------------------------------------
%
%  (end of template credits)
%

\chapter{Conclusion}
\label{conclusion}

After reading this tutorial it should be easy to quickly write great building automation applications and if that gets tiring, try sombrero(\url{http://grill.github.com/sombrero/}), a web interface to control EIB/KNX enabled houses. Sombrero is the first application to use Scalimero. 

%
% EOF
%


%\printindex
\end{document}

%
% EOF
%
